\subsection{\textbf{Sprint 0}}

O projeto teve início com uma reunião com o cliente, onde foi apresentada a ideia do projeto, como surgiu e as inspirações que foram consideradas para realizá-lo, além de terem sido esclarecidas as dúvidas que foram surgindo. Juntamente com o Lucas, os \ac{ages} IV do projeto, Rafael Cardoso e Vitor Hugo, montaram \textit{sketches} de cada tela, alterando conforme as ideias do time iam surgindo e sendo validadas pelo próprio stakeholder.

Após a reunião de apresentação do projeto, o time se juntou para definir as tecnologias que seriam utilizadas no desenvolvimento. Em seguida, dividimo-nos em squads de quatro integrantes, com o objetivo de gerar diferentes ideias para as telas do aplicativo e ter mais opções. Cada integrante, especialmente os \ac{ages} I, criou telas de acordo com sua compreensão dos requisitos do stakeholder.

Após a criação das primeiras ideias para o aplicativo, o time se reuniu novamente para selecionar os principais pontos positivos de cada \textit{mockup}, a fim de evitar desperdício de trabalho.

Com uma base sólida do que deveria ser feito, os gerentes do projeto decidiram dividir as telas entre as squads. Meu squad foi responsável pelos objetivos e também fui encarregado de iniciar o \textit{style guide} do projeto. O squad organizou um horário para discutir suas ideias, a fim de chegar a uma tela final. Como não sabia quais componentes os outros squads poderiam precisar, criei um \textit{style guide} mais abrangente, incluindo alguns componentes específicos utilizados nas telas de objetivos.

Nas últimas aulas da sprint, como o squad havia concluído as telas, cada um ajudou os demais conforme possível. Como eu já tinha conhecimento em animações e \textit{flows} no Figma, trabalhei em conjunto com o Squad do Fogo, auxiliando principalmente na criação de componentes que mudam de cor ao serem clicados, bem como no botão \textit{switch}, utilizados nas telas de objetivos e na tela principal.